\begin{abstract}
We propose and formalize a \emph{Computational Action Principle} (CAP) in the unified Holographic Polar Arithmetic (HPA) and $\Omega$-theory framework \cite{Ma2025HPA,Ma2025FoP,Ma2025HPOT}. The guiding thesis is operational: physical ``laws'' are not externally imposed rules but minimal-cost error-correction and steady-state constraints required for finite-resolution readout to remain self-consistent over long horizons. The core tension is layered. At the ontic layer, the universe is modeled by a continuous unitary scan and local quantum updates (Weyl pairs, PQCA). At the readout layer, observers access only discrete projections with finite information capacity, inducing canonical tick structures (Ostrowski/Zeckendorf) and structural mismatch measured by discrepancy.

CAP upgrades this mismatch into a variational principle: dynamics selects configurations minimizing a \emph{least-discrepancy} functional subject to local covariance and implementation constraints. Under standard closure assumptions (locality, diffeomorphism invariance, and second-order metric equations), the macroscopic gravitational field equation is forced to be the Einstein equation (with a cosmological constant term) by Lovelock-type uniqueness; discrepancy and implementation costs enter only through an effective stress tensor and potential.

We further motivate an $\Omega$ action in which the Fisher-information amplitude of an information density field is minimally coupled to gravity, and the \emph{routing overhead} of compiling local updates to nearest-neighbor circuits appears as a computational lapse field, providing an operational interpretation of gravitational time dilation. Gauge fields arise as compensating connections for local phase readout errors implied by Weyl complementarity, while matter is modeled as topologically locked phase defects.

Finally, we provide reproducible toy experiments: numerical star-discrepancy comparisons across scan slopes illustrate the finite-depth optimality of the golden branch, and an FFT-based Poisson solver verifies that a localized defect source produces an approximate $1/r$ phase potential and an inverse-square ``phase pressure'' acceleration in the near-field regime.
\end{abstract}

\textbf{Keywords}: Holographic Polar Arithmetic (HPA); $\Omega$ theory; Computational Action Principle; least discrepancy; Fisher information; Weyl pair; routing overhead; emergent gravity; gauge connection; topological defect.

\paragraph{Conventions.}
Unless otherwise stated, $\log$ denotes the natural logarithm. ``mod $1$'' refers to reduction in $\RR/\ZZ$. We use the mostly-plus metric signature $(-,+,+,+)$ and set $c=1$ in theoretical derivations unless explicitly restored.

\tableofcontents
\newpage


